\documentclass[UTF8,aspectratio=169]{ctexbeamer}
\usetheme{Madrid}
\usecolortheme{default}
\usepackage{amsmath}
\usepackage{booktabs}
\usepackage{siunitx}
\usepackage{graphicx}

\title{Distillation + Raw CDF Refinement\\(v2 Engineered-Feature Architecture)}
\author{Mie Spray Postprocessing Project}
\date{\today}

\begin{document}

\begin{frame}
  \titlepage
\end{frame}

\begin{frame}{Message}
\textbf{Claim:} a Stage-2 heteroscedastic teacher trained on fitted curves can be refined toward raw CDF observations using knowledge distillation, with regime-dependent weighting that respects the time-dependent disappearance of usable raw data.

\vspace{0.4em}
\textbf{What this deck covers:}
\begin{itemize}
  \item why direct raw-series supervision alone is insufficient
  \item how CDF coverage regimes partition the time axis
  \item the distillation loss design: raw NLL + KD KL + onset + anchor + shape priors
  \item v2-specific changes: $A$-scaled feature space and physical-space loss
  \item the student model architecture and training procedure
\end{itemize}
\end{frame}

\begin{frame}{Motivation: Why Not Just Train on Raw Series?}
\textbf{The problem with raw CDF alone:}
\begin{itemize}
  \item raw CDF observations thin out with time --- by $t\sim\SI{3}{ms}$ many operating conditions have lost $>70\%$ of their measurements
  \item fitting only on available raw data produces models that collapse toward zero or become erratic in the late-time regime
  \item the fitted-curve teacher (Stage-2 NLL) is well-defined everywhere on $[0, 5]\,\si{ms}$ but was never calibrated against raw observations
\end{itemize}

\vspace{0.3em}
\textbf{Solution:} use the teacher's Gaussian predictions as a regularizing prior in time regions where raw data are sparse, and use raw CDF as direct supervision where coverage is strong.
\end{frame}

\begin{frame}{CDF Coverage Regime Labeling}
Time is binned into $\Delta t=\SI{0.1}{ms}$ intervals over $[0, 5]\,\si{ms}$. For each operating-condition group, smoothed raw CDF sample counts determine three regimes:

\vspace{0.3em}
\begin{description}
  \item[\texttt{raw\_reliable}] smoothed coverage $\geq70\%$ of peak --- raw CDF is the primary supervision source
  \item[\texttt{raw\_uncertain}] coverage has dropped below $70\%$ but above $20\%$ --- transitional zone
  \item[\texttt{teacher\_only}] coverage below $20\%$ or fewer than 4 raw samples per bin --- teacher Gaussian becomes the only supervision
\end{description}

\vspace{0.3em}
\textbf{Implementation:} \texttt{build\_time\_bin\_regimes} produces a per-group, per-bin regime table. Transitions require $\geq2$ consecutive bins to prevent isolated dips from triggering premature regime changes.
\end{frame}

\begin{frame}{Regime-Dependent Supervision Weights}
Each grid point receives two independent weight channels:

\vspace{0.3em}
\begin{center}
\begin{tabular}{lcc}
\toprule
Regime & Raw weight $w_{\mathrm{raw}}$ & KD weight $w_{\mathrm{kd}}$ \\
\midrule
\texttt{raw\_reliable}   & 1.0 & 0.25 \\
\texttt{raw\_uncertain}  & 0.0 & 1.0  \\
\texttt{teacher\_only}   & 0.0 & 0.75 \\
\bottomrule
\end{tabular}
\end{center}

\vspace{0.3em}
\textbf{Design rationale:}
\begin{itemize}
  \item in \texttt{raw\_reliable}: trust raw data fully, keep mild KD regularization
  \item in \texttt{raw\_uncertain}: drop raw supervision entirely, let teacher KL dominate
  \item in \texttt{teacher\_only}: slightly reduce KD weight to acknowledge teacher extrapolation uncertainty
\end{itemize}

Both channels are further modulated by a global time-bin rebalancing weight (inverse-frequency, clipped to $[0.5, 2.0]$).
\end{frame}

\begin{frame}{v2 Adaptation: $A$-Scaled Feature and Loss Space}
\textbf{Recall from the v2 engineered-feature deck:}
\[
A = \left(\frac{\Delta P}{\rho_a}\right)^{1/4} d_n^{1/2}, \qquad
\hat\mu(t) = \mu(t)/A, \qquad
\hat\sigma(t) = \sigma(t)/A
\]

\vspace{0.3em}
\textbf{What changes in distillation:}
\begin{itemize}
  \item features are built via \texttt{build\_feature\_matrix\_np}, which returns \texttt{(features, a\_scale, canonical)} --- the dataset now carries \texttt{a\_scale} per sample
  \item teacher forward pass: $\mu_{\mathrm{teacher}}^{\mathrm{phys}} = A\,\hat\mu_{\mathrm{teacher}}$
  \item student forward pass: $\mu_{\mathrm{student}}^{\mathrm{phys}} = A\,\hat\mu_{\mathrm{student}}$
  \item \textbf{all losses are computed in physical space (mm)}: raw NLL, KD KL, anchor penalty, monotonicity and concavity priors
\end{itemize}

\vspace{0.3em}
This ensures the loss magnitude is comparable across operating conditions with different $A$ values, and raw CDF targets (which are always in mm) can be used directly.
\end{frame}

\begin{frame}{Student Model Architecture}
\textbf{Output dimension:} 3 (vs.\ teacher's 2)
\[
\text{student output} = [\hat\mu,\;\log\hat\sigma^2,\;\text{onset\_logit}]
\]

\textbf{Initialization:}
\begin{itemize}
  \item all shared layers are copied from the teacher checkpoint
  \item final linear layer: first 2 output rows copied from teacher; onset row initialized to zero (sigmoid output $\to 0.5$)
\end{itemize}

\vspace{0.3em}
\textbf{Feature set:} identical to the teacher --- \texttt{a\_only} (5 features) or \texttt{a\_plus\_log\_a} (6 features). The student is a pointwise MLP, but batching is sequence-wise so derivative priors operate on contiguous trajectories.
\end{frame}

\begin{frame}{Composite Loss Function}
\[
\mathcal{L} = \mathcal{L}_{\mathrm{raw}} + \mathcal{L}_{\mathrm{KD}} + \lambda_{\mathrm{onset}}\,\mathcal{L}_{\mathrm{onset}} + \lambda_{\mathrm{anchor}}\,\mathcal{L}_{\mathrm{anchor}} + \lambda_{d_1}\,P_{d_1} + \lambda_{d_2}\,P_{d_2}
\]

\vspace{0.2em}
\begin{description}
  \item[$\mathcal{L}_{\mathrm{raw}}$] Gaussian NLL on physical $\mu$, $\sigma^2$ vs.\ raw CDF, weighted by $w_{\mathrm{raw}}$
  \item[$\mathcal{L}_{\mathrm{KD}}$] forward KL divergence $\mathrm{KL}(q_{\mathrm{student}}\|p_{\mathrm{teacher}})$ in physical space, weighted by $w_{\mathrm{kd}}$ and a teacher confidence factor $\min(1, \sigma_{\mathrm{ref}}/\sigma_{\mathrm{teacher}})$
  \item[$\mathcal{L}_{\mathrm{onset}}$] BCE on the onset logit vs.\ a ramp target $\min(t/t_{\mathrm{ramp}}, 1)$, active only for $t\leq\SI{0.2}{ms}$
  \item[$\mathcal{L}_{\mathrm{anchor}}$] quadratic penalty on $\mu(t\approx0)$, ramping down over the first $\SI{0.15}{ms}$
  \item[$P_{d_1}$] monotonicity prior: $\mathrm{ReLU}(-\Delta\mu)^2$
  \item[$P_{d_2}$] concavity prior: $\mathrm{ReLU}(+\Delta^2\mu)^2$, gated to activate after $t\geq\SI{0.5}{ms}$
\end{description}
\end{frame}

\begin{frame}{Teacher Confidence Weighting}
The KD loss includes a per-point confidence weight:
\[
w_{\mathrm{conf}}(t) = \mathrm{clamp}\!\left(\frac{\sigma_{\mathrm{ref}}}{\sigma_{\mathrm{teacher}}(t)},\;0.25,\;1.0\right), \qquad \sigma_{\mathrm{ref}} = \SI{10}{mm}
\]

\textbf{Purpose:} where the teacher's predicted standard deviation is large (i.e.\ the teacher itself is uncertain), the KD signal is down-weighted. This prevents the student from being pulled toward poorly constrained teacher predictions in extrapolation regions.

\vspace{0.3em}
The effective KD weight at each grid point is $w_{\mathrm{kd}} \times w_{\mathrm{conf}} \times w_{\mathrm{time\_bin}}$.
\end{frame}

\begin{frame}{Training Procedure}
\textbf{Data split:} random $70/15/15$ train/val/test on CDF wide-format samples.

\vspace{0.3em}
\textbf{Sampling:} \texttt{WeightedRandomSampler} with inverse onset-bin frequency, so early-onset and late-onset samples are seen equally often.

\vspace{0.3em}
\textbf{Optimization:}
\begin{itemize}
  \item AdamW, $\mathrm{lr}=3\times10^{-3}$, weight decay from teacher config
  \item gradient clipping at norm $1.0$
  \item early stopping on validation loss, patience $=20$
\end{itemize}

\vspace{0.3em}
\textbf{Sanity gates before full training:}
\begin{enumerate}
  \item static check: one batch forward pass, verify finite loss and valid teacher confidence bounds
  \item tiny overfit check: 30 steps on $\leq8$ samples, confirm loss decreases and onset logit moves toward zero near $t=0$
\end{enumerate}
\end{frame}

\begin{frame}{Script Usage}
\textbf{CLI invocation:}
\begin{verbatim}
python distillation_plus_raw_series.py \
    stage2_engineered_nll_a_only_20260410_020607 \
    --epochs 150 --patience 20
\end{verbatim}

\vspace{0.3em}
\textbf{Key flags:}
\begin{itemize}
  \item positional \texttt{run\_dir}: Stage-2 NLL run directory (teacher)
  \item \texttt{--no-train}: run all diagnostics and checks without starting the training loop
  \item \texttt{--series-split \{clean, all\}}: which raw series to load
  \item \texttt{--device \{auto, cpu, cuda\}}
\end{itemize}

\vspace{0.3em}
\textbf{Outputs saved to \texttt{runs\_mlp/distill\_cdf\_onset\_v2\_<timestamp>/}:}
\begin{itemize}
  \item \texttt{best\_model\_refinement.pt}, \texttt{refine\_config.json}
  \item \texttt{epoch\_loss.csv}, \texttt{test\_summary.csv}, \texttt{loss\_curves.png}
  \item \texttt{teacher\_run\_dir.txt} for provenance
\end{itemize}
\end{frame}

\begin{frame}{What Still Needs Evaluation}
\begin{itemize}
  \item compare student vs.\ teacher on held-out raw CDF in the \texttt{raw\_uncertain} regime --- does the student better track the raw scatter?
  \item run the sparse-feature stability diagnostics on the student: the $A$-scaling should still prevent oscillatory $d_n$ and $P_{\mathrm{inj}}$ sweeps
  \item evaluate onset prediction accuracy vs.\ the engineered hydraulic delay
  \item compare \texttt{a\_only} vs.\ \texttt{a\_plus\_log\_a} teacher $\to$ student distillation
  \item longer-term: wire the onset head into the inference path as a soft masking mechanism
\end{itemize}
\end{frame}

\end{document}
